
\documentclass[a4paper]{article}
\usepackage[pdftex]{graphicx}
\usepackage[utf8]{inputenc}
\usepackage{enumerate}
\usepackage{siunitx}
\sisetup{locale=DE} 
\usepackage{href-ul}
\hypersetup{
	colorlinks=true,
	linkcolor=blue,
	urlcolor=blue}
\usepackage{icomma}
\usepackage{amssymb}
\usepackage{geometry}
\geometry{a4paper, top=15mm, left=15mm, right=15mm, bottom=15mm,
	headsep=10mm, footskip=12mm}

\begin{document}
%Titel
{\bf Übungsblatt Wellen}\\
\begin{enumerate}[1.]
%Aufgabe 1
\begin{minipage}{0.5\textwidth}
\item Zwei Lautsprecher befinden sich in einem x-y-Koordinatensystem an den Punkten $L_1(0|80)$ und $L_2(0|-80)$. Mit einem druckempfindlichen Mikrophon wird das in der Umgebung der Lautsprecher entstehende Wellenfeld untersucht. Die Membranen der beiden Lautsprecher schwingen gleichzeitig mit der Frequenz $f=\SI{680}{\hertz}$. Das Mikrophon wird vom Punkt $A(1,50|0)$ aus längs der Geraden $g$ verschoben. Dabei werden Intensitätsmaxima und Intensitätsminima der Druckschwankungen registriert. 
\end{minipage}
\hspace{0.5 cm}
\begin{minipage}{0.5\textwidth}
\includegraphics[width=6 cm]{lauti237}
\end{minipage}

\begin{enumerate}[a)]
\item Erklären Sie ausführlich das Zustandekommen eines Intensitätsmaximums.
\item Im Punkt $B(1,50|1,30)$ registriert man ein Intensitätsmaximum 2. Ordnung. Bestätigen Sie, dass die Schallge\-schwindigkeit den Betrag $c=\SI{340}{\meter\over \second}$ hat. 
\item Berechnen Sie die Anzahl derjenigen Punkte auf der Geraden $g$, in denen ein Intensitätsmaximum nachgewiesen werden kann.
\end{enumerate}

\item In einem zweiten Versuch ist nur noch ein Lautsprecher in Betrieb. Dieser Lautsprecher befindet sich im Ursprung des x-y-Koordinatensystems. Die Membran des Lautsprechers schwingt mir der Frequenz $f=\SI{680}{\hertz}$. Das Mikrophon wird mit einer konstanten Geschwindigkeit 
$\vec{v}_M$ vom Punkt $A(1,50|0)$ aus längs der x-Achse nach rechts bewegt. Dabei registriert das Mikrophon Druckschwankungen mit der konstanten Frequenz $f^*$.

\begin{enumerate}[a)]
\item Leiten Sie eine Formel her, mit der die Frequenz $f^*$ aus $c,~f$ und dem Betrag $v_M$ der Geschwindigkeit $\vec{v}_M$ berechnet werden kann. Erläutern Sie dabei Ihre Überlegungen.
\item Berechnen Sie $f^*$ für $v_M=\SI{2,00}{\meter\over \second}$.
\item Das Mikrophon wird nun mit einer konstanten Geschwindigkeit $\vec{v}$ vom Punkt 
$A(1,50|0)$ aus auf der Geraden $g$ in Richtung zunehmender y-Werte bewegt. Begründen Sie, dass das Mikrophon bei diesem Versuch Druckschwankungen registriert, deren Frequenz mit wachsender Entfernung des Mikrophons vom Lautsprecher abnimmt.
\end{enumerate} 

\end{enumerate} 
\fbox{
	\begin{minipage}{0.5\textwidth}
		Weitere Arbeitsblätter gibt es unter 
		
		\href{https://www.okuyakl.de}{www.okuyakl.de}
	\end{minipage}
	\hfill
	\begin{minipage}{0.4\textwidth}
		\includegraphics[width=1.5 cm]{../../viecher/zwe03}
		\includegraphics[width=2 cm]{../../viecher/afanticon1}
		
\end{minipage}}
%Ende
\end{document}